\documentclass[10pt,a4paper]{amsart}
\usepackage{amsmath,amssymb,amsthm,graphicx}
% Retained preamble for the Stacks Project chapter snapshot in this directory.
% The chapter's own preamble.tex is a copy of the shared Stacks Project preamble; the block
% below is sliced MECHANICALLY, line for line, from it. Nothing was added, removed or
% redefined; the wrapper class, the xr-hyper/hyperref wiring and the externaldocument list
% are omitted because the assembler supplies the document class.
%
% lines taken from preamble.tex: 3, 5, 9, 11, 12, 156, 157, 158, ... (51 lines total)
%
% package coverage restored from the chapter's own preamble (2026-10-13): xy, xy-2cell

\usepackage[all]{xy}
\xyoption{2cell}
\UseAllTwocells

\newenvironment{reference}{\comment}{\endcomment}
\newenvironment{slogan}{\comment}{\endcomment}
\newenvironment{history}{\comment}{\endcomment}
\theoremstyle{plain}
\newtheorem{theorem}[subsection]{Theorem}
\newtheorem{proposition}[subsection]{Proposition}
\newtheorem{lemma}[subsection]{Lemma}
\theoremstyle{definition}
\newtheorem{definition}[subsection]{Definition}
\newtheorem{example}[subsection]{Example}
\newtheorem{exercise}[subsection]{Exercise}
\newtheorem{situation}[subsection]{Situation}
\theoremstyle{remark}
\newtheorem{remark}[subsection]{Remark}
\newtheorem{remarks}[subsection]{Remarks}
\numberwithin{equation}{subsection}
\def\lim{\mathop{\mathrm{lim}}\nolimits}
\def\colim{\mathop{\mathrm{colim}}\nolimits}
\def\Spec{\mathop{\mathrm{Spec}}}
\def\Hom{\mathop{\mathrm{Hom}}\nolimits}
\def\Ext{\mathop{\mathrm{Ext}}\nolimits}
\def\SheafHom{\mathop{\mathcal{H}\!\mathit{om}}\nolimits}
\def\SheafExt{\mathop{\mathcal{E}\!\mathit{xt}}\nolimits}
\def\Sch{\mathit{Sch}}
\def\Mor{\mathop{\mathrm{Mor}}\nolimits}
\def\Ob{\mathop{\mathrm{Ob}}\nolimits}
\def\Sh{\mathop{\mathit{Sh}}\nolimits}
\def\NL{\mathop{N\!L}\nolimits}
\def\CH{\mathop{\mathrm{CH}}\nolimits}
\def\proetale{{pro\text{-}\acute{e}tale}}
\def\etale{{\acute{e}tale}}
\def\QCoh{\mathit{QCoh}}
\def\Ker{\mathop{\mathrm{Ker}}}
\def\Im{\mathop{\mathrm{Im}}}
\def\Coker{\mathop{\mathrm{Coker}}}
\def\Coim{\mathop{\mathrm{Coim}}}
\DeclareMathSymbol{\boxtimes}{\mathbin}{AMSa}{"02}
\def\QCohstack{\mathcal{QC}\!\mathit{oh}}
\def\Cohstack{\mathcal{C}\!\mathit{oh}}
\def\Spacesstack{\mathcal{S}\!\mathit{paces}}
\def\Quotfunctor{\mathrm{Quot}}
\def\Hilbfunctor{\mathrm{Hilb}}
\def\Curvesstack{\mathcal{C}\!\mathit{urves}}
\def\Polarizedstack{\mathcal{P}\!\mathit{olarized}}
\def\Complexesstack{\mathcal{C}\!\mathit{omplexes}}
\def\Pic{\mathop{\mathrm{Pic}}\nolimits}
\def\Picardstack{\mathcal{P}\!\mathit{ic}}
\def\Picardfunctor{\mathrm{Pic}}
\def\Deformationcategory{\mathcal{D}\!\mathit{ef}}
\usepackage{cleveref}
\crefname{originallemma}{original lemma}{original lemmas}
\Crefname{originallemma}{Original lemma}{Original lemmas}
\makeatletter
\newlabel{divisors-lemma-finite-locally-free-has-norm}{{0BD2}{source}{}{sourcepointer-1}{}}
\newlabel{divisors-lemma-finite-locally-free-has-norm@cref}{{[originallemma][0][]0BD2}{[1][0][]source}{}{}{}}
\makeatother
\crefname{originallemma}{original lemma}{original lemmas}
\Crefname{originallemma}{Original lemma}{Original lemmas}
\makeatletter
\newlabel{section-discriminant}{{0BVH}{source}{}{sourcepointer-2}{}}
\newlabel{section-discriminant@cref}{{[originallemma][0][]0BVH}{[1][0][]source}{}{}{}}
\makeatother
\title{Stacks, Discriminants: lemma-norm-different-in-discriminant}
\author{The Stacks Project Authors}
\date{stacks-project:discriminant.tex:local-snapshot-20261012; GNU FDL 1.2 or later}
\begin{document}
\maketitle
\noindent\textit{Editorial note.} Statement span is lines 1862-1867 of sources/stacks-discriminant/pinned-original-source.tex; proof span is lines 1869-1904, the complete author proof (34 body lines).\par\medskip
\noindent\textit{Editorial note.} Navigation stubs: 2 cited-but-not-included labels are supplied as non-load-bearing pointers with their pinned tags where the snapshot tag index has them.\par\medskip
\noindent\textit{Editorial note.} extra\_packages is empty on purpose: the retained preamble already loads xy with options and a second declaration clashes in the assembler wrapper.\par\medskip
\noindent\textit{Editorial note.} Text is retained verbatim, including the author's own typos and placeholders.\par\medskip
\noindent\textit{Editorial note.} Queue index 20 of queue\_discriminant.json (pb=1295, stmt 1862-1867, proof 1869-1904).\par\medskip
\noindent\textit{Allowed standard prerequisites (editorial).}\begin{itemize}
\item Explicit prerequisite boundary. The cited inputs are prerequisites; the argument in the retained span is not replaced by any cited result.
\end{itemize}\medskip
\ifdefined\hypertarget\hypertarget{sourcepointer-1}{}\fi
\noindent\textit{Source reference.} Original lemma 0BD2: sources/stacks-discriminant/pinned-original-source.tex, source lines 1904--1904. Cited by the retained statement or proof but not part of the collected spans; supplied only as a navigation pointer. The argument does not depend on its content.\par\medskip
\ifdefined\hypertarget\hypertarget{sourcepointer-2}{}\fi
\noindent\textit{Source reference.} Original lemma 0BVH: sources/stacks-discriminant/pinned-original-source.tex, source lines 694--759. Cited by the retained statement or proof but not part of the collected spans; supplied only as a navigation pointer. The argument does not depend on its content.\par\medskip
\begin{lemma}
\label{lemma-norm-different-in-discriminant}
Let $f : Y \to X$ be a finite flat morphism of Noetherian schemes.
Then $\text{Norm}_f : f_*\mathcal{O}_Y \to \mathcal{O}_X$ maps
$f_*\mathfrak{D}_f$ into the ideal sheaf of the discriminant $D_f$.
\end{lemma}

\begin{proof}
The norm map is constructed in
Divisors, Lemma \ref{divisors-lemma-finite-locally-free-has-norm}
and the discriminant of $f$ in Section \ref{section-discriminant}.
The question is affine local, hence we may assume $X = \Spec(A)$,
$Y = \Spec(B)$ and $f$ given by a finite locally free ring map $A \to B$.
Localizing further we may assume $B$ is finite free as an $A$-module.
Choose a basis $b_1, \ldots, b_n \in B$ for $B$ as an $A$-module.
Denote $b_1^\vee, \ldots, b_n^\vee$ the dual basis of
$\omega_{B/A} = \Hom_A(B, A)$ as an $A$-module.
Since the norm of $b$ is the determinant of $b : B \to B$ as an
$A$-linear map, we see that
$\text{Norm}_{B/A}(b) = \det(b_i^\vee(bb_j))$.
The discriminant is the principal closed subscheme of $\Spec(A)$
defined by $\det(\text{Trace}_{B/A}(b_ib_j))$.
If $b \in \mathfrak{D}_{B/A}$ then
there exist $c_i \in B$ such that
$b \cdot b_i^\vee = c_i \cdot \text{Trace}_{B/A}$ where
we use a dot to indicate the $B$-module structure on $\omega_{B/A}$.
Write $c_i = \sum a_{il} b_l$.
We have
\begin{align*}
\text{Norm}_{B/A}(b)
& =
\det(b_i^\vee(bb_j)) \\
& =
\det( (b \cdot b_i^\vee)(b_j)) \\
& =
\det((c_i \cdot \text{Trace}_{B/A})(b_j)) \\
& =
\det(\text{Trace}_{B/A}(c_ib_j)) \\
& =
\det(a_{il}) \det(\text{Trace}_{B/A}(b_l b_j))
\end{align*}
which proves the lemma.
\end{proof}

\end{document}
